\par\noindent\textbf{(1)} 展开 \((a+a^\dagger)^4\) 并使用 \([a,a^\dagger]=1\) 正规序，可得 \(\langle n|(a+a^\dagger)^4|n\rangle=6n^2+6n+3\)。因此 \(\langle n|x^4|n\rangle=3(2n^2+2n+1)[\hbar/(2m\omega)]^2\)。
\par\noindent\textbf{(2)} 非简并微扰给 \(E_n^{(1)}=\lambda\langle n|x^4|n\rangle=3\lambda\hbar^2(2n^2+2n+1)/(4m^2\omega^2)\)。因此相邻能级差的修正 \(E_{n+1}^{(1)}-E_n^{(1)}=3\lambda\hbar^2(n+1)/(m^2\omega^2)\)，能级间隔随激发数增大。
\par\noindent\textbf{(3)} 对 \(n=0\)，\(E_0^{(1)}=3\lambda\hbar^2/(4m^2\omega^2)\)。由于 \(x^4\) 为偶宇称，只能连接偶数态；例如 \(\langle2|x^4|0\rangle=6\sqrt2[\hbar/(2m\omega)]^2\)，基态矢量的 \(|2\rangle\) 一阶系数为该矩阵元乘 \(\lambda/(E_0^{(0)}-E_2^{(0)})\)。
\par\noindent\textbf{(4)} 衡量小扰动的基态无量纲参数为 \(\epsilon=\lambda\hbar/(m^2\omega^3)\ll1\)。在高 \(n\) 态，四次能量修正约按 \(\lambda n^2\) 增长而无扰能量约按 \(\hbar\omega n\) 增长，即使基态满足微扰小量，大 \(n\) 区也会失效。
